\clearpage
\subsection{Sprint 2}

A \textit{Sprint} 2 teve como objetivo dar continuidade à primeira versão navegável do ReadRace e conectar funcionalidades que ainda estavam parcialmente isoladas. Uma das prioridades foi concluir o fluxo de desafios entre amigos, pendente da \textit{Sprint} 1. Também estava prevista a integração do registro de leitura ao \ac{XP} acumulado, ao nível do usuário, à sequência de dias, ao progresso dos desafios e às conquistas. Essa integração era importante para que registrar uma leitura produzisse efeitos perceptíveis no restante da experiência do aplicativo.

O planejamento incluía as telas de Perfil e Minhas Conquistas, a possibilidade de curtir e descurtir posts e a correção de pendências de navegação, apresentação visual, acessibilidade e testes automatizados. Inicialmente, também seriam trabalhadas a participação em comunidades, a página da comunidade e a adição de livros. Entretanto, a redução da duração da \textit{sprint} de quatro para duas semanas exigiu rever os compromissos. A apresentação de entrega da \textit{Sprint} 2 registra essa revisão de 21 tarefas inicialmente apresentadas para 18 no plano final. As três tarefas retiradas foram transferidas para a \textit{Sprint} 3, delimitando o que seria efetivamente entregue naquele intervalo.

Além dessa revisão das tarefas da \textit{sprint}, conversamos com os clientes sobre a retirada dos eventos globais do escopo do produto. A proposta já estava muito ampla, e esse alinhamento permitiu concentrar o trabalho no que seria viável entregar. Essa negociação foi importante para que a redução de escopo fosse uma decisão compartilhada e para manter uma expectativa clara sobre o resultado do semestre.

Na organização da equipe, passei a acompanhar a frente de Social, biblioteca e qualidade, junto ao colega Guilherme Bohm, também \ac{AGES} IV. Criei um grupo de WhatsApp para facilitar a comunicação e compartilhei as tarefas e seus links: curtir e descurtir posts (\href{https://github.com/ReadRace-AGES/ReadRace/issues/101}{issue \#101}), corrigir o retorno do detalhe do livro para a Busca (\href{https://github.com/ReadRace-AGES/ReadRace/issues/102}{\#102}), remover rotas de teste fora do produto (\href{https://github.com/ReadRace-AGES/ReadRace/issues/103}{\#103}) e substituir o aviso local da biblioteca pelo Toast compartilhado (\href{https://github.com/ReadRace-AGES/ReadRace/issues/104}{\#104}). Deixei os colegas escolherem as tarefas e me coloquei, junto ao Guilherme, à disposição para esclarecer dúvidas de requisitos e de desenvolvimento.

O acompanhamento ocorreu pelo WhatsApp e diretamente nas aulas; com os \ac{AGES} III do \textit{squad}, utilizei mais o Discord. A divisão funcionou muito bem para manter um trabalho mais consistente e direto com os colegas. Conseguimos reduzir a necessidade de os \ac{AGES} IV assumirem implementações para compensar dificuldades da equipe, preservando espaço para que os demais integrantes conduzissem suas entregas. Os registros do repositório mostram essa distribuição: as curtidas foram implementadas no \href{https://github.com/ReadRace-AGES/ReadRace/pull/112}{PR \#112}, de \texttt{arthurgreco}; a retirada de rotas, no \href{https://github.com/ReadRace-AGES/ReadRace/pull/114}{PR \#114}, de \texttt{gugs123}; e a troca do aviso pelo Toast, no \href{https://github.com/ReadRace-AGES/ReadRace/pull/123}{PR \#123}, de \texttt{oEduardoCastro}.

Minha contribuição direta de implementação nessa \textit{sprint} foi a correção de navegação da issue \#102, integrada pelo \href{https://github.com/ReadRace-AGES/ReadRace/pull/124}{PR \#124}. Ao abrir um livro pela Busca, o retorno podia levar à aba errada porque Busca e Meus Livros compartilhavam o mesmo caminho de rota. Separei os caminhos por aba, mantendo a tela de detalhe compartilhada, de forma a preservar o termo e os resultados da pesquisa e a pilha de navegação de origem. Também foram adicionados testes de retorno nos estados de carregamento, erro e sucesso e após o registro de progresso. Essa tarefa mostrou como um problema de navegação aparentemente pequeno podia comprometer a continuidade da experiência de uso.

Ao final, a equipe concluiu as 18 tarefas do escopo revisado, incluindo seis novas telas. Nos desafios entre amigos, foram entregues a criação de desafios por páginas ou por livro, a escolha da obra, o acompanhamento do placar e a apresentação dos resultados de vitória, derrota ou empate. Esses fluxos foram construídos em contribuições de diferentes colegas, como os \href{https://github.com/ReadRace-AGES/ReadRace/pull/115}{PRs \#115}, \href{https://github.com/ReadRace-AGES/ReadRace/pull/121}{\#121}, \href{https://github.com/ReadRace-AGES/ReadRace/pull/122}{\#122} e \href{https://github.com/ReadRace-AGES/ReadRace/pull/126}{\#126}. Minha atuação de gestão apoiava esse esforço coletivo por meio da comunicação e do acompanhamento das tarefas.

Também foram entregues Perfil, Minhas Conquistas e as interações de curtida e descurtida. Houve correções na navegação e na exibição de capas, ajustes de layout no celular e a inclusão de testes do \textit{mobile} na integração contínua. Como entrega adicional da equipe, foi estruturada a infraestrutura \ac{AWS} com configuração de \textit{deploy} automático. A integração da leitura ao total de \ac{XP} e ao nível foi registrada no \href{https://github.com/ReadRace-AGES/ReadRace/pull/127}{PR \#127}; a sequência de dias, no \href{https://github.com/ReadRace-AGES/ReadRace/pull/125}{PR \#125}; e conquistas e seu desbloqueio automático, no \href{https://github.com/ReadRace-AGES/ReadRace/pull/118}{PR \#118}. Essas mudanças conectaram o registro de leitura à progressão do usuário.

Essa etapa deu continuidade ao trabalho com registro de leitura e cálculo de \ac{XP} no qual eu havia atuado na \textit{Sprint} 1. Para mim, a evolução do produto tornou mais evidente que uma regra aparentemente localizada pode afetar diversas partes da aplicação. O progresso de leitura precisava manter coerência com nível, desafios e conquistas, o que aumentava a importância de alinhar o comportamento esperado entre as diferentes frentes de desenvolvimento.

O cronograma situou a apresentação da \textit{Sprint} 2 aos \textit{stakeholders} e o planejamento da \textit{Sprint} 3 em 2 de outubro de 2026, às 19h15. Para a mesma data, às 21h, estavam previstas a retrospectiva e a entrega do relatório de andamento no Fluxo \ac{AGES}. Conduzi a apresentação dessa \textit{sprint}, assumindo de forma mais direta um desafio que havia estabelecido para a \ac{AGES} IV: melhorar minhas habilidades de comunicação e apresentação. Preparar o conteúdo e explicar as entregas aos clientes passou a ser uma parte importante da minha atuação de gestão.

A principal lição de processo foi a necessidade de ajustar o escopo à duração real da \textit{sprint} e explicitar as transferências de tarefas. A conclusão do escopo revisado precisava ser compreendida junto dessa mudança, mantendo visibilidade sobre o que continuava pendente. Para a etapa seguinte, permaneceram demandas de autenticação e disponibilização do aplicativo instalável, além da expansão dos clubes, comunidades, biblioteca e recursos de personalização. Assim, a \textit{Sprint} 2 ampliou a integração do produto, mas ainda não representou o encerramento das funcionalidades planejadas para o semestre.
